% Generated from validated synthesis. Do not hand-edit.
\section*{Retrospective Signals}
\addcontentsline{toc}{section}{Retrospective Signals}
\smallskip
\noindent\textbf{能力を増やすだけでなく、振る舞いを設計する層が独立した} InstructGPT、Chain-of-Thought、FLAN、Self-Instructは、post-training、推論時のprompting、instruction tuning、synthetic instruction dataという異なるレバーを可視化した。 \hfill{\footnotesize p.~\pageref{pkg:steering-reasoning}}\par
\smallskip
\noindent\textbf{scaleはparameter数からcompute・access・runtimeの問題へ広がった} ChinchillaとPaLM、OPT、BLOOM、GLM-130B、FlashAttentionを並べると、training allocation、model distribution、memory/IO efficiencyが別々の制約として見える。 \hfill{\footnotesize p.~\pageref{pkg:scale-access-systems}}\par
\smallskip
\noindent\textbf{生成メディアはmodel単体からdata・personalization・distributionの生態系へ進んだ} Flamingo、DALL-E 2、Stable Diffusionを中心に、LAION-5BやDreamBooth、音声・動画・三次元の生成研究が、multimodal理解とmedia generationの設計空間を複数方向へ広げた。 \hfill{\footnotesize p.~\pageref{pkg:multimodal-media}}\par
\smallskip
\noindent\textbf{生成結果がcode、environment action、retrievalへ接続された} Copilot、ReAct、VPT、Atlasなどにより、model outputを外部systemへ渡すinterface layerが、model本体とは別の技術領域として前景化した。 \hfill{\footnotesize p.~\pageref{pkg:code-action-retrieval}}\par
\smallskip
\noindent\textbf{対話UIの普及と同時にevaluationとcontrolが前景化した} ChatGPT research previewを年末の統合点として、LaMDA、Whisper、helpful/harmless RLHF、Sparrow、Constitutional AI、HELMまでを追うと、利用者向けinterfaceと制御・評価の層が結び付き始めたことが分かる。 \hfill{\footnotesize p.~\pageref{pkg:dialogue-interface-control} / p.~\pageref{pkg:evaluation-layer} / p.~\pageref{pkg:annual-synthesis}}\par
\medskip
\begin{claimboundary}[Retrospective scope]
本号は2022年を後日確認可能になった一次情報も用いて再構成するRetrospective Specialである。Coverage Periodと制作時点を同一視せず、vendor / project / author claimの境界は各記事とTechnical Notesで明示する。
\end{claimboundary}
\medskip
\tableofcontents
